\section{Parametric Study}\label{sec:results}

\subsection{Elastic Modulus Sensitivity}

The $n=2$ flexural mode frequency depends strongly on the wall elastic
modulus $E$.  Table~\ref{tab:E_sweep} shows the predicted frequency for
a free shell across the physiologically relevant range.

\begin{table}[htbp]
  \centering
  \caption{Flexural mode $n=2$ frequency vs.\ elastic modulus
  ($a = \SI{0.18}{\metre}$, $c = \SI{0.12}{\metre}$, $h = \SI{10}{\milli\metre}$).}
  \label{tab:E_sweep}
  \begin{tabular}{rcc}
    \hline
    $E$ (MPa) & $f_2$ (Hz) & In ISO~2631 range? \\
    \hline
    0.05 & 3.6 & Marginal \\
    0.10 & 4.4 & Yes \\
    0.20 & 5.6 & Yes \\
    0.50 & 8.3 & Yes \\
    1.00 & 11.5 & No \\
    2.00 & 16.1 & No \\
    \hline
  \end{tabular}
\end{table}

The \SIrange{4}{8}{\hertz} ISO~2631 abdominal resonance range requires
$E < \SI{0.5}{\mega\pascal}$, corresponding to relaxed, passive musculature.
Active muscle bracing increases $E$ by a factor of 10--100, moving the
resonance well above the infrasound range.

\subsection{Multi-Parameter Sensitivity}

A full factorial analysis over 486 parameter combinations (6 values of $E$,
3 semi-major axes, 3 aspect ratios, 3 wall thicknesses, 3 loss tangents)
reveals that 37\% of configurations place the $n=2$ mode in the
\SIrange{5}{10}{\hertz} range and 41\% in the broader \SIrange{4}{8}{\hertz}
ISO~2631 range.  The dominant parameter is $E$, followed by geometry and
intra-abdominal pressure.

\subsection{Multi-Layer Wall Model}

The abdominal wall comprises approximately five distinct tissue layers:
peritoneum, transversalis fascia, internal oblique muscle, external oblique
muscle, subcutaneous fat, and skin.  Using classical laminate theory, the
effective composite modulus for a relaxed wall is
$E_\mathrm{eff} \approx \SI{0.075}{\mega\pascal}$
($h_\mathrm{total} = \SI{25.1}{\milli\metre}$), giving
$f_2 = \SI{4.6}{\hertz}$---within 4\% of the homogeneous model prediction.
Muscle tension is confirmed as the dominant factor: tensed muscle
($E_\mathrm{eff} = \SI{2.7}{\mega\pascal}$) shifts $f_2$ to
\SI{23.1}{\hertz}.

\subsection{Boundary Condition Effects}

The free-shell model provides a lower bound on frequency. Partial constraints
(spine, pelvis, ribs) raise the $n=2$ frequency by an estimated factor of
1.3--3.0$\times$, depending on the angular extent of the constrained boundary.
Even with a $2\times$ constraint multiplier, the mode remains within the
ISO~2631 range for soft tissue.

\subsection{ISO~2631 Validation}

Table~\ref{tab:iso} compares model predictions against published abdominal
resonance data.

\begin{table}[htbp]
  \centering
  \caption{Model predictions vs.\ published abdominal resonance frequencies.}
  \label{tab:iso}
  \begin{tabular}{lcc}
    \hline
    Body type & $f_2$ (Hz) & ISO~2631 range? \\
    \hline
    Lean ($E = 0.2$, $a = 15$ cm) & 6.3 & Yes \\
    Average ($E = 0.1$, $a = 18$ cm) & 4.0 & Yes \\
    Large ($E = 0.08$, $a = 20$ cm) & 3.1 & Marginal \\
    \hline
  \end{tabular}
\end{table}

